\chapter{Diagnosi e verifica delle modifiche}
\label{chap:debug}

Questo capitolo serve a localizzare un errore e a verificare una modifica. I
comandi ordinari di esecuzione e valutazione restano nel
Capitolo~\ref{chap:esecuzione}; qui sono riportati soltanto i controlli utili
quando un risultato non è quello atteso.

\section{Diagnosi di un nodo}

La regola principale è cercare il primo artefatto errato: detection,
riconciliazione delle istanze, geometria, archi e infine diametri. Per conservare
lo stato intermedio, abilitare temporaneamente i dump in \filepath{.env}:

\begin{lstlisting}[language={}]
DEBUG_OUTPUT_ENABLED=true
DEBUG_OUTPUT_DIR=./output/debug
\end{lstlisting}

Rieseguire quindi un solo nodo:

\begin{lstlisting}[language={}]
python run.py 7
\end{lstlisting}

I file interni vengono scritti sotto
\filepath{<DEBUG_OUTPUT_DIR>/node7/}. Non fanno parte della consegna e gli
strumenti che li leggono non sono eseguiti automaticamente da
\filepath{run.py}. Al termine riportare \envvar{DEBUG_OUTPUT_ENABLED} a
\code{false}.

\begin{contracttable}
Componente assente & \filepath{<WORK_DIR>/detection_cache/} & Verificare che il codice
compaia nella cache della relativa immagine. Se manca, il problema precede la
ReID ed è da cercare tra YOLO, crop e decoder. \\

Componente duplicato &
\filepath{<WORK_DIR>/multiview/node<id>/scene_nodes.json} & Confrontare le
immagini di supporto delle istanze con lo stesso codice. Supporti complementari
indicano una possibile mancata fusione. \\

Separazione ReID errata &
\filepath{<DEBUG_OUTPUT_DIR>/node<id>/reid_geometry.json} & Controllare
posizione 3D, appartenenza dei crop e vincoli tra osservazioni prima di
modificare il clustering. \\
\midrule

Arco mancante o errato &
\filepath{<WORK_DIR>/multiview/node<id>/} & Verificare prima gli
\code{SceneNode}, quindi \filepath{vggt_geometry.json} e infine gli archi in
\filepath{fused_graph.json}. \\

Diametro incoerente &
\filepath{<DEBUG_OUTPUT_DIR>/node<id>/diameters.json} & Confrontare
osservazioni, associazione iniziale e propagazione. I codici \code{250} e
\code{252} non hanno un valore millimetrico nella mappa corrente. \\
\end{contracttable}

I comandi diagnostici più utili sono:

\begin{lstlisting}[language={}]
python tools/debug/debug_crops.py 7
python tools/debug/analyze_reid_false_positives.py 7
python tools/debug/plot_reid_geometry.py 7 15 95
python tools/debug/debug_graph_projection.py 7
\end{lstlisting}

Il primo crea una tavola dei crop per \code{SceneNode}; il secondo, se sono
disponibili label e annotazioni, confronta i cluster con il ground truth; il
terzo visualizza la separazione 3D per i codici richiesti; l'ultimo proietta
nodi e archi sulle immagini e salva un
\filepath{consistency_report.json}. Le distanze 2D mostrate da quest'ultimo sono
un aiuto visivo e non un'evidenza usata per costruire il grafo.

\manualnote[ID composti]{Tra i comandi elencati,
\filepath{debug_crops.py} e \filepath{debug_graph_projection.py} convertono
l'argomento in numero intero e non accettano identificativi come \code{15_1}.
Per questi nodi occorre adattare il parsing oppure richiamare la funzione
interna passando l'ID come stringa.}

\section{Verifica delle modifiche}

Una modifica alla pipeline va confrontata sugli stessi nodi e sugli stessi
input. Quando cambiano modello o dataset è preferibile usare una nuova
\envvar{WORK_DIR}, così cache e risultati precedenti restano disponibili.

La verifica minima consiste nel:

\begin{enumerate}
  \item rieseguire almeno un nodo che attraversi il ramo modificato;
  \item confrontare il primo artefatto interessato e il
  \filepath{fused_graph.json} finale;
  \item eseguire valutazione e visualizzazione come descritto nel
  Capitolo~\ref{chap:esecuzione};
  \item conservare comando, configurazione e differenze osservate.
\end{enumerate}

L'edge F1 di \filepath{evaluate.py} usa il miglior abbinamento biunivoco trovato
tra istanze predette e annotate, tenendo conto anche della struttura. Va quindi
letto come misura della qualità del grafo sotto quel mapping, non come verifica
della stabilità degli identificativi prodotti.

Il repository non contiene una suite di test dedicata. Prima di rifattorizzare
ReID, costruzione degli archi o diametri è opportuno congelare almeno un caso
per ciascuno di questi comportamenti: fusione multi-vista con VGGT, fallback di
connessione e propagazione del diametro.

\section{Strumenti specialistici}

\filepath{evaluation/standard_comparison/} mantiene il confronto strutturale
richiamato da \filepath{evaluate.py}. \filepath{shurikode_symbols.py} gestisce
simboli e porte condivisi con QGIS, mentre \filepath{package_qgis_geojson.py}
prepara gli archivi. Gli script in \filepath{tools/reporting/} e
\filepath{tools/examples/} sono specifici o dimostrativi: non fanno parte di
\filepath{run.py} e non sono test.
